\chapter{Requerimientos del sistema}

El presente capítulo describe los requerimientos que el sistema debe satisfacer para
constituir una solución válida al problema de navegación asistida para personas con
discapacidad visual. Estos requerimientos se derivan de las necesidades del usuario final,
de las restricciones del entorno de operación y de las características de la arquitectura
modular definida en el capítulo anterior. Sirven como base directa para el capítulo de
análisis, en el cual se justifica la selección de tecnologías a partir de los criterios
aquí establecidos.

Los requerimientos se organizan en dos categorías: funcionales, que describen qué debe
hacer el sistema, y no funcionales, que definen las restricciones de calidad bajo las
cuales debe operar.

% ==============================================================================
\section{Requerimientos Funcionales}
% ==============================================================================

Los requerimientos funcionales se agrupan conforme a la arquitectura modular del sistema,
siguiendo la misma organización de módulos descrita en el marco metodológico.

% ------------------------------------------------------------------------------
\subsection*{RF -- Módulos de interacción}
% ------------------------------------------------------------------------------

\begin{itemize}

    \item[\textbf{RF-01}] \textbf{Captura continua de voz con detección de actividad.}
    El sistema debe capturar la entrada de audio del usuario de forma continua mediante el
    micrófono del dispositivo y detectar automáticamente los segmentos en los que hay señal
    de voz activa, descartando los periodos de silencio para optimizar el procesamiento.

    \item[\textbf{RF-02}] \textbf{Filtrado y normalización de audio.}
    El sistema debe aplicar supresión de ruido ambiental y normalización de nivel sobre la
    señal de audio antes de entregarla al módulo de procesamiento de lenguaje natural, con
    el fin de reducir errores de reconocimiento en entornos con ruido de fondo.

    \item[\textbf{RF-03}] \textbf{Reconocimiento de voz e interpretación de intención.}
    El sistema debe transcribir la señal de voz a texto, clasificar la intención del usuario
    ---indicar un destino, solicitar descripción del entorno, controlar la sesión--- y
    extraer las entidades relevantes para ejecutar la acción correspondiente.

    \item[\textbf{RF-04}] \textbf{Síntesis y reproducción de instrucciones por voz.}
    El sistema debe convertir en audio las instrucciones de navegación, las alertas y las
    respuestas generadas por el sistema, reproduciéndolas al usuario mediante síntesis de
    voz con lenguaje natural, frases cortas y referencias espaciales claras.

    \item[\textbf{RF-05}] \textbf{Control de sesión por voz.}
    El sistema debe reconocer comandos de voz para iniciar, pausar, reanudar y detener la
    sesión de navegación, sin requerir ningún tipo de interacción física o visual por parte
    del usuario.

    \item[\textbf{RF-06}] \textbf{Exposición completa de funciones a través de voz.}
    El módulo de interacción debe exponer la totalidad de las funciones del sistema mediante
    comandos de voz, de forma que ninguna función quede accesible exclusivamente por medios
    táctiles o visuales. Esto incluye iniciar y detener la navegación, indicar destinos,
    solicitar descripciones del entorno y pedir rutas alternativas.

\end{itemize}

% ------------------------------------------------------------------------------
\subsection*{RF -- Módulos de percepción}
% ------------------------------------------------------------------------------

\begin{itemize}

    \item[\textbf{RF-07}] \textbf{Captura continua de imágenes y preprocesamiento.}
    El sistema debe adquirir fotogramas de forma continua mediante la cámara del dispositivo
    y aplicar operaciones de preprocesamiento ---redimensionamiento, normalización e
    igualación de contraste--- para estandarizar la entrada al módulo de visión.

    \item[\textbf{RF-08}] \textbf{Detección de objetos con localización espacial.}
    El módulo de visión debe detectar los objetos relevantes para la navegación presentes en
    cada fotograma ---incluyendo al menos personas, obstáculos estáticos, puertas, escaleras
    y desniveles--- generando para cada uno una región de interés con su nivel de confianza
    asociado.

    \item[\textbf{RF-09}] \textbf{Clasificación de objetos y estimación de nivel de riesgo.}
    El sistema debe asignar una categoría semántica a cada objeto detectado y determinar su
    nivel de riesgo para la trayectoria actual, distinguiendo entre elementos que representan
    un peligro inmediato y elementos que forman parte del entorno navegable sin representar
    una amenaza.

    \item[\textbf{RF-10}] \textbf{Estimación de distancia a objetos detectados.}
    El sistema debe estimar la distancia entre el usuario y cada objeto detectado a partir
    de la información visual disponible y de la representación espacial del entorno, con el
    fin de determinar la urgencia de la alerta correspondiente.

    \item[\textbf{RF-11}] \textbf{Validación cruzada con sensores de proximidad.}
    El sistema debe contrastar la información visual con las mediciones de los sensores de
    proximidad del dispositivo para confirmar o descartar la presencia de obstáculos,
    reduciendo la tasa de falsas alarmas antes de notificar al módulo de navegación.

    \item[\textbf{RF-12}] \textbf{Generación y emisión de alerta de obstáculo.}
    Al confirmar la presencia de un obstáculo relevante, el sistema debe notificar al usuario
    por voz de forma inmediata, especificando la naturaleza del obstáculo y su posición
    relativa, y enviar las coordenadas y el nivel de riesgo al módulo de navegación para
    desencadenar el recálculo de ruta si corresponde.

    \item[\textbf{RF-13}] \textbf{Inferencia de relaciones contextuales entre objetos.}
    El sistema debe inferir relaciones espaciales y contextuales entre los elementos
    detectados ---como el estado de una puerta o el grado de obstrucción de un pasillo---
    para enriquecer la representación del entorno más allá de la detección individual
    de objetos.

\end{itemize}

% ------------------------------------------------------------------------------
\subsection*{RF -- Módulos de representación}
% ------------------------------------------------------------------------------

\begin{itemize}

    \item[\textbf{RF-14}] \textbf{Estimación continua de la posición del usuario.}
    El sistema debe mantener y actualizar de forma continua una estimación de la posición y
    orientación del usuario dentro del entorno, integrando el movimiento detectado con la
    información perceptual disponible.

    \item[\textbf{RF-15}] \textbf{Construcción incremental del mapa del entorno.}
    El sistema debe construir y actualizar progresivamente un mapa del entorno a medida que
    el usuario se desplaza, sin requerir un mapa previo del espacio como condición para
    operar.

    \item[\textbf{RF-16}] \textbf{Representación de zonas libres y ocupadas.}
    El mapa generado debe distinguir explícitamente entre zonas navegables y zonas ocupadas
    por obstáculos, proporcionando al módulo de planeación una estructura de datos adecuada
    para calcular trayectorias seguras.

    \item[\textbf{RF-17}] \textbf{Actualización dinámica del mapa.}
    El mapa del entorno debe actualizarse en respuesta a los cambios detectados, reflejando
    la aparición o desaparición de obstáculos de forma que la representación espacial
    corresponda siempre al estado actual del entorno. La frecuencia mínima de esta
    actualización se establece en RNF-03.

\end{itemize}

% ------------------------------------------------------------------------------
\subsection*{RF -- Módulo de navegación y rutas}
% ------------------------------------------------------------------------------

\begin{itemize}

    \item[\textbf{RF-18}] \textbf{Cálculo de ruta óptima local.}
    El módulo de navegación debe calcular una trayectoria que conecte la posición actual del
    usuario con su destino, minimizando el costo del recorrido y garantizando que la ruta
    propuesta evite todas las zonas ocupadas registradas en el mapa.

    \item[\textbf{RF-19}] \textbf{Actualización periódica de la ruta durante la navegación.}
    El sistema debe recalcular y actualizar la ruta activa de forma periódica durante una
    sesión de navegación, incorporando los cambios detectados en el mapa para mantener la
    trayectoria vigente coherente con el estado actual del entorno. La frecuencia mínima
    de este recálculo se establece en RNF-04.

    \item[\textbf{RF-20}] \textbf{Recálculo de ruta ante obstáculos confirmados.}
    Cuando el módulo de percepción confirme la presencia de un obstáculo en la trayectoria
    activa, el módulo de navegación debe generar una ruta alternativa de forma inmediata,
    sin esperar al ciclo periódico de actualización establecido en RNF-04.

    \item[\textbf{RF-21}] \textbf{Generación de instrucciones de movimiento.}
    El módulo de dirección debe traducir cada punto de la ruta calculada en una instrucción
    de movimiento discreta ---avanzar, girar a la derecha, girar a la izquierda,
    detenerse--- expresada en términos de distancia y ángulo comprensibles, y entregarla
    al módulo de salida para su emisión por voz.

    \item[\textbf{RF-22}] \textbf{Resolución de destinos sobre el mapa del entorno.}
    El sistema debe permitir que el usuario indique como destino un elemento reconocido y
    registrado en el mapa ---por ejemplo, una puerta o la salida--- o un punto explícito,
    y guiarlo hacia él mediante la navegación local, sin depender de servicios de
    cartografía externos.

    \item[\textbf{RF-23}] \textbf{Solicitud de ruta alternativa por voz.}
    El usuario debe poder solicitar por voz una ruta alternativa en cualquier momento de la
    navegación. El sistema debe analizar su viabilidad y aplicarla, o notificar al usuario
    si no existe una alternativa disponible.

\end{itemize}

% ==============================================================================
\section{Requerimientos No Funcionales}
% ==============================================================================

Los requerimientos no funcionales establecen los criterios de calidad bajo los cuales el
sistema debe operar. Son determinantes en la selección de tecnologías, ya que imponen
restricciones concretas sobre el rendimiento, la precisión, la portabilidad y la seguridad
de cada módulo.

% ------------------------------------------------------------------------------
\subsection*{RNF -- Rendimiento}
% ------------------------------------------------------------------------------

\begin{itemize}

    \item[\textbf{RNF-01}] \textbf{Latencia de detección y alerta de obstáculos.}
    El sistema debe detectar, clasificar y emitir la alerta de un obstáculo en un tiempo
    máximo de 2 segundos desde la captura del fotograma en el que éste aparece, para
    garantizar que la notificación al usuario sea útil antes de que la situación represente
    un riesgo.

    \item[\textbf{RNF-02}] \textbf{Latencia de reconocimiento de voz.}
    El sistema debe procesar e interpretar un comando de voz en menos de 500~ms desde que
    el usuario termina de hablar, de forma que la respuesta o acción resultante no interrumpa
    el flujo natural de la interacción.

    \item[\textbf{RNF-03}] \textbf{Frecuencia de actualización del mapa.}
    La representación espacial del entorno debe actualizarse a una frecuencia mínima de
    10~Hz, de modo que los cambios en el entorno queden reflejados en el mapa sin un desfase
    perceptible entre la posición real del usuario y su posición estimada en el mapa.

    \item[\textbf{RNF-04}] \textbf{Frecuencia de actualización de la ruta activa.}
    La ruta de navegación debe recalcularse al menos una vez cada diez segundos durante una
    sesión activa, garantizando que la trayectoria incorpore los cambios recientes del
    entorno sin imponer una carga computacional continua al sistema. Este valor constituye
    la restricción temporal asociada al comportamiento funcional descrito en RF-19.

    \item[\textbf{RNF-05}] \textbf{Eficiencia computacional en hardware limitado.}
    El sistema debe poder ejecutarse en hardware de propósito general con recursos limitados,
    como los disponibles en dispositivos móviles o sistemas embebidos de bajo costo ---por
    ejemplo, una Raspberry Pi 4 o equivalente--- cumpliendo los tiempos de respuesta
    establecidos sin requerir infraestructura de cómputo en la nube para su operación básica.

\end{itemize}

% ------------------------------------------------------------------------------
\subsection*{RNF -- Precisión y confiabilidad}
% ------------------------------------------------------------------------------

\begin{itemize}

    \item[\textbf{RNF-06}] \textbf{Precisión en la detección de obstáculos.}
    El módulo de visión debe alcanzar una tasa de detección correcta mínima del 90\% bajo
    condiciones de iluminación normales, manteniendo la tasa de falsas alarmas por debajo
    del 5\% para evitar interrupciones innecesarias durante la navegación.

    \item[\textbf{RNF-07}] \textbf{Tasa de reconocimiento correcto de comandos de voz.}
    El módulo de procesamiento de lenguaje natural debe interpretar correctamente al menos
    el 92\% de los comandos emitidos en condiciones de ruido ambiental de hasta 60~dB.

    \item[\textbf{RNF-08}] \textbf{Disponibilidad y degradación controlada.}
    El sistema debe permanecer operativo durante toda la sesión de navegación sin reinicios
    no solicitados. Ante el fallo de un módulo no crítico, debe degradarse de forma
    controlada, notificando al usuario la pérdida de funcionalidad y manteniendo operativas
    las funciones restantes.

    \item[\textbf{RNF-09}] \textbf{Comportamiento conservador ante incertidumbre.}
    Ante cualquier situación de baja confianza en la información del entorno ---como
    detección con umbral bajo o pérdida de señal de posicionamiento--- el sistema debe
    optar por la acción más conservadora: advertir al usuario y sugerir detenerse hasta
    recuperar información confiable. Es preferible una alerta de precaución adicional que
    omitir un obstáculo real.

\end{itemize}

% ------------------------------------------------------------------------------
\subsection*{RNF -- Usabilidad y accesibilidad}
% ------------------------------------------------------------------------------

\begin{itemize}

    \item[\textbf{RNF-10}] \textbf{Interfaz exclusivamente por voz.}
    El sistema debe poder operarse en su totalidad mediante comandos de voz, sin requerir
    interacción táctil o visual por parte del usuario durante la navegación. Esta restricción
    aplica a la experiencia completa del usuario y complementa lo establecido en RF-06,
    que define la responsabilidad funcional del módulo de interacción.

    \item[\textbf{RNF-11}] \textbf{Claridad y comprensibilidad de las instrucciones.}
    Las instrucciones generadas por el módulo de salida deben utilizar lenguaje natural,
    frases cortas y referencias espaciales concretas ---izquierda, derecha, metros--- de
    forma que sean comprensibles sin ambigüedad para un usuario sin conocimientos técnicos
    y sin entrenamiento previo en el uso del sistema.

    \item[\textbf{RNF-12}] \textbf{Retroalimentación proactiva ante eventos relevantes.}
    El sistema debe emitir retroalimentación auditiva de forma proactiva ante eventos
    significativos ---detección de un obstáculo, cambio de ruta, pérdida de señal--- sin
    esperar a que el usuario lo solicite explícitamente.

\end{itemize}

% ------------------------------------------------------------------------------
\subsection*{RNF -- Portabilidad y mantenibilidad}
% ------------------------------------------------------------------------------

\begin{itemize}

    \item[\textbf{RNF-13}] \textbf{Operación autónoma sin conexión a internet.}
    Los módulos de percepción, representación, planeación local e interacción deben
    funcionar de forma autónoma sin requerir conexión a internet. Los componentes que
    aprovechen la conexión para mejorar su calidad deben contar con una alternativa local
    que se active automáticamente en ausencia de conexión, manteniendo operativas todas
    las funciones del sistema.

    \item[\textbf{RNF-14}] \textbf{Modularidad e independencia entre componentes.}
    Los módulos de voz, visión, representación y navegación deben estar desacoplados e
    intercomunicados mediante interfaces bien definidas, de modo que cualquiera de ellos
    pueda ser modificado o sustituido sin afectar el funcionamiento de los demás, siempre
    que se respeten las interfaces establecidas.

    \item[\textbf{RNF-15}] \textbf{Facilidad de actualización de modelos.}
    La arquitectura debe permitir sustituir los modelos de percepción o inferencia por
    versiones mejoradas sin requerir cambios estructurales en el resto del sistema, de forma
    que el desempeño pueda actualizarse conforme avance el estado del arte en las tecnologías
    empleadas.

\end{itemize}